\section*{Limitations}
\label{app:limitations}

Four of five generators are open-weight; GPT-5.4 is a non-reproducible Azure preview snapshot (\texttt{gpt-5.4}, Appendix~\ref{app:judges}) included to confirm that frontier outputs do not escape the benchmark's difficulty gradients. All primary claims replicate on the four open-weight generators independently of it.

We adopt a deterministic programmatic scorer, rather than a human or LLM annotator, as the reference signal, since it is the only reference that scales to $321{,}648$ evaluations without reintroducing the judge-reliability question under study. A 120-record human study (Appendix~\ref{app:prog-judge-validation}) validates this choice at $92.5$--$98.3\%$ agreement on three of four metrics, but only $82.5\%$ on parameter structure, where the scorer penalises schema-valid extra argument keys that annotators accept. Each of the 120 records was scored by a single annotator, so we cannot report an inter-annotator agreement statistic for the human labels themselves; the $92.7\%$ scorer-validation figure and the human-verdict judge rankings (including the QwQ-32B 1st$\to$4th reversal) should be read as agreement with one calibrated annotator's judgement, not with a consensus reference. A multi-annotator replication with an inter-annotator agreement check is left to future work. The parameter-structure gap does not drive our findings: excluding parameter structure from the aggregate ranking changes it by $\leq0.2$\,pp (Table~\ref{tab:param-robustness}), and a maximally generous upper-bound correction moves every judge's alignment by $\leq0.09$\,pp with the ranking unchanged (Table~\ref{tab:eq2-bound}). These checks establish that the six-judge \emph{ranking} is robust to this gap; they do not establish that the \emph{absolute} alignment percentages reported throughout the paper are unaffected; by construction every headline alignment number incorporates the parameter-structure metric at its current, imperfectly-validated scoring rule. We report both programmatic and human-verdict rankings wherever a best-judge claim is made, since they diverge on some records; a schema-aware rescoring of the corpus is left to future work.

The scorer also does not model tool redundancy: the synthesized tool inventory is not guaranteed to expose a single canonical path for a given subtask, so a generator that reaches an equivalent result via a structurally different tool (e.g., a dedicated aggregation tool vs.\ composing two simpler ones) can be penalised by $p^{\text{tool}}$ and $p^{\text{seq}}$ even though the substitution is conceptually correct. The 120-record human study already surfaces a version of this: $9$ of $21$ tool-selection disagreements between the programmatic scorer and human annotators involve semantically equivalent tools with different names, or defensible extra calls the scorer treats as out-of-reference (Appendix~\ref{app:prog-judge-validation}). We have not separately quantified how much of the corpus admits such redundant pathways, nor built a schema-aware tool-equivalence correction analogous to the parameter-structure bound in Table~\ref{tab:eq2-bound}; both are left to future work.

Records are synthetically generated rather than mined from enterprise traces, since a verified reference for every record is required by our paired GT/no-GT protocol and is not obtainable at this scale from proprietary systems. A two-level quality gate and the same human study ($92.7\%$ agreement) validate the pipeline, but domain drift from real deployments remains an open validity concern; a real-trace comparison is blocked on corpus access and left to future work. All prompts are scoped to single-shot, stateless planning, and headline findings are comparative rather than absolute-magnitude, which limits, but does not eliminate, exposure to this gap.

Our six main judges are all general-purpose models. To test generalisation to judge-specialised models, we added Prometheus-2~\citep{kim2024prometheus2} across all five generators and both GT conditions (Table~\ref{tab:full-results}, row Prom2). It does not cluster with the general-purpose judges under either condition ($20$--$30$\,pp lower alignment, near-chance pairwise agreement; Table~\ref{tab:interjudge}), with the gap concentrated on parameter structure and query coverage (Table~\ref{tab:e4-permetric}). It also shows a negative GT lift opposite to most judges, with higher programmatic alignment without GT than with it, which does not contradict its own published reference-free result against \emph{human} judgement~\citep[Appendix F]{kim2024prometheus2}, a different comparison axis. It is excluded from all six-judge statistics and does not affect any headline finding; a broader judge-specialised roster is left to future work.

GPT-5.4 simultaneously serves as generator, judge, and rewrite meta-judge. Its self-bias on its own generations falls within its range on other generators (Appendix~\ref{app:self-preference}), except for sequence accuracy, where it over-credits itself ($+0.172$ vs.\ $\leq+0.008$ elsewhere); we find no aggregate self-preference but cannot rule out this metric-localised effect. We treat this as an open finding rather than a resolved one: the check is a post-hoc correlational analysis of the existing evaluation grid, not a controlled ablation that isolates GPT-5.4's role (Appendix~\ref{app:self-preference}), and the one metric where a self-bias signal is measurable is also the one we do not resolve. This is compounded by GPT-5.4 being a non-reproducible Azure preview snapshot with no fixed version string (Appendix~\ref{app:judges}): unlike the effects we bound elsewhere in this section, this specific signal cannot be independently re-verified on the same model in future work.

The RQ4 (temperature) and RQ6 (prompt format) ablations were each extended to a second (judge, generator) pairing. Temperature insensitivity replicates more tightly; the prompt-format advantage replicates only partially and reverses on hard queries, so we treat it as judge/generator-dependent rather than a uniformly dominant lever. A full grid across all judges and generators is left to future work.

Difficulty rewriting is unanimously validated on $58.1\%$ of easy$\to$medium and $93.9\%$ of medium$\to$hard records; the latter is our primary difficulty-degradation evidence, and the "three difficulty tiers" framing used throughout the paper should be read with medium as a robustness check rather than a fully independent, calibrated tier.

The without-GT ceiling is partly attributable to the prompt's default-to-$1.0$ rubric: the C2 ablation (Table~\ref{tab:c2-ceiling}) shows the ceiling shifts by only ${\leq}{+1.0}$\,pp under an alternative $0.5$-default for the three strongest generators, but by $+4.1$ and $+5.6$\,pp for the two weakest (Llama-3.1-8B, SmolLM3-3B). Task difficulty remains the primary driver of the ceiling for capable generators, but its exact height is not fully capacity-independent; the ceiling is somewhat prompt-dependent specifically on weaker generators, a nuance the abstract's ``structural ceiling'' framing compresses. Headline figures use unweighted topology averages over an intentionally imbalanced corpus; count-weighted averaging shifts figures by $\leq1.5$\,pp with no ranking change (Appendix~\ref{app:pertopology-table}).

We study judge reliability at evaluation time and do not test whether the failure modes we identify would carry over if the same judges were used as a training signal, e.g., a reward model in RLHF/DPO fine-tuning of a tool-calling agent, or as an automatic gate for model selection. This distinction matters because the mechanisms are different in kind, not just degree: an evaluation-time error changes a single reported number, while a training-time error changes the objective an agent is optimized against. Three of our findings would each shift the risk profile differently. GT-exposure over-anchoring (Finding 2) would only affect training pipelines that include the reference trace in the reward model's input, and would push the agent toward mimicking the surface form of provided demonstrations rather than being penalised for missing steps, since our case studies show the with-GT judge crediting an incomplete plan whenever it superficially resembles the reference (Appendix~\ref{app:anchoring-cases}). The without-GT verdict-compression ceiling (\S\ref{sec:exp:ceiling}) is the higher-risk case for training use, since a reward model built on such a judge would supply a near-constant reward on hard queries regardless of the agent's actual correctness, weakening the gradient signal exactly where the agent most needs correction. The correlated failure across judges (six-judge ensemble matching, not exceeding, the best individual judge; \S\ref{sec:exp:ceiling}) implies that this weak-gradient risk would not be mitigated by ensembling several of our evaluated judges into one reward signal, since they fail on the same records for the same structural reason rather than independently. Testing any of this directly would require training runs with judge-derived rewards under controlled bias conditions, which is outside this paper's scope; we flag it as a concrete, motivated direction for future work rather than a speculative one.
